% Source: https://artofproblemsolving.com/community/c3h1813326_word_problems
% Stable identity: aops-topic-1813326-q01
Let the car from $A$ have speed $v$ km per hour, so the car from $B$ has speed $v+10$. In the first scenario they close the gap at $2v+10$ km per hour for $5$ hours, so
\[
AB=5(2v+10)=10v+50 .
\]
In the second scenario the car from $B$ drives $150$ km, taking $\dfrac{150}{v+10}$ hours, while the car from $A$ drives for $4.5$ hours longer. Together they cover $AB$:
\[
v\left(\dfrac{150}{v+10}+4.5\right)+150=10v+50\iff \dfrac{150v}{v+10}=5.5v-100 .
\]
Multiplying by $2(v+10)$: $300v=(11v-200)(v+10)=11v^2-90v-2000$, so
\[
11v^2-390v-2000=0\iff(v-40)(11v+50)=0 .
\]
Since $v>0$, $v=40$, and $AB=10\cdot40+50=\boxed{450\text{ km}}$.

Check: in the second scenario the car from $B$ (at $50$ km per hour) drives $150$ km in $3$ hours, while the car from $A$ drives $7.5$ hours at $40$ km per hour, i.e.\ $300$ km; and $300+150=450$.
